\begin{proof}
\begin{enumerate}[label=\arabic*.]
\item \textbf{The actual cusp marking and component.}
Put $N=M_0-I$. Since $N^2=0$ and $\operatorname{rank}N=3$,
$\operatorname{im}N=\ker N$ over $\mathbb R$.
Choose a basis $u_1,u_2,u_3$ of this space and vectors $v_j$
with $Nv_j=u_j$. Applying $N$ to a relation among these six
vectors first kills its $v$ coefficients, and then independence
of the $u_j$ kills the others. Thus
\[
 P=(u_1,u_2,u_3,v_1,v_2,v_3),\qquad
 N_s=\begin{pmatrix}0&I_3\\0&0\end{pmatrix}
\]
satisfy $NP=PN_s$ and $M_0P=P(I+N_s)$.
We retain $\det P$, including its sign. The chosen period
component in the original coordinates is exactly $PD_3$;
the map $F\mapsto PF$ is its holomorphic Grassmannian
identification with the positive logarithmic component $D_3$.
Every $\rho(\gamma)$ preserves $PD_3$, because
$P^{-1}\rho(\gamma)P$ has positive real determinant.

For a primitive integral cusp, the two rank-three subgroups
$\operatorname{im}N$ and $\ker N$ of $\Lambda$ have the same
rational span and are primitive, hence equal. Choose integral
$u_j$ and integral preimages $v_j$. The induced map
$N:\Lambda/\ker N\longrightarrow\operatorname{im}N$ is an
isomorphism, taking $[v_j]$ to $u_j$. Every lattice vector is
therefore an integral combination of the six chosen vectors.
This proves $P\in\mathrm{GL}(6,\mathbb Z)$ and supplies the
marking of \ref{maps:src:F.2.1} with its actual basis map.

\item \textbf{The logarithmic bundle extension.}
Let $B=\mathbb P^1\setminus\{p_1,\ldots,p_k\}$, which is
open because $k\ge1$. On the punctured coordinate disc at
$p_0$ put $s=\log t/(2\pi i)$. In the given local flat
bundle the equivariance is $(s,z)\mapsto(s+1,M_0z)$.
Define
\[
 T(s)=P(I+sN_s),\qquad
 T(s)^{-1}=(I-sN_s)P^{-1}.
\]
Then $T(s+1)=M_0T(s)$. Consequently the coordinate
$c=T(s)^{-1}z$ is single-valued. Glue a trivial bundle across
$t=0$ to the actual flat bundle by $c\mapsto T(s)c$.
This gives a holomorphic bundle $V_D\to B$ with the specified
flat monodromy on $B^*$.
In the marked coordinates the untwisting is
$I-sN_s=\exp(-sN_s)$, since $N_s^2=0$.

The central Grassmannian chart is
$F_C=\operatorname{im}\binom C{I_3}$,
$C\in\operatorname{Mat}_3(\mathbb C)$. Undoing the gauge gives
\[
 T(s)F_C=P\operatorname{im}\binom{sI_3+C}{I_3}.
\]
The logarithmic coefficient and the correction matrix in
\eqref{maps:eq:F.O6} are therefore exactly those in the
specified cusp marking.

\item \textbf{An open target containing the central graph chart.}
Over $B^*$ take the transverse-plane bundle in the component
$PD_3$. At $p_0$ adjoin every plane in the graph chart $F_C$,
which is precisely the set transverse to $\ker(N_s)_\mathbb C$
in the extended marking. Denote this subset of
$\operatorname{Gr}_3(V_D)$ by $\mathscr U$.
For $t\ne0$, its membership in this chart is
\[
 \det\bigl(y(t)I_3+\operatorname{Im}C\bigr)>0,
 \qquad y(t)=-\frac{\log|t|}{2\pi}.
 \tag{F.O7}\label{maps:eq:F.O7}
\]
Indeed, for $Z=sI_3+C$, subtraction of the two column blocks
and exchange of their three-row blocks gives
\[
 \det\left(\binom Z{I_3},\binom{\overline Z}{I_3}\right)
       =(2i)^3\det(\operatorname{Im}Z).
\]
The positive component is the one containing $Z=iyI_3$.
For a bounded neighbourhood of any fixed $C_0$, sufficiently
small $|t|$ gives $y>\|\operatorname{Im}C\|$ uniformly.
The Neumann inverse applies to
$I_3+y^{-1}\operatorname{Im}C$, and it also applies with
$\operatorname{Im}C$ replaced by $r\operatorname{Im}C$ for
$0\le r\le1$. The determinant consequently retains the
positive sign of $y^3$. A full neighbourhood of $(0,C_0)$
thus lies in $\mathscr U$, proving openness.

Its projection onto $B$ is surjective. In local Grassmannian
coordinates it is $(t,C)\mapsto t$, with differential
$(\delta t,\delta C)\mapsto\delta t$ and right inverse
$\delta t\mapsto(\delta t,0)$, so it is a holomorphic
submersion. Its fibre at $p_0$ is the displayed $\mathbb C^9$
chart. Over $B^*$ it is the holomorphic bundle with fibre
$PD_3$, identified by $P$ with the LSAP fibre of
\ref{maps:src:F.3.8}.

\item \textbf{A continuous section retaining every prescribed jet.}
For the prescribed matrix coefficients $A_0,\ldots,A_r$,
put $C_*(t)=\sum_{j=0}^r A_jt^j$. Choose a small closed disc
$K=\{|t|\le\epsilon\}$, with $\epsilon\le1$, such that
\[
 -\frac{\log\epsilon}{2\pi}>
                  \sum_{j=0}^r\|A_j\|+1.
\]
Its graph is a holomorphic section of $\mathscr U$ near $K$.
The complement $B\setminus K$ is connected and noncompact,
so $K$ is Runge.

We give the fibre topology needed for the relative extension.
A transverse plane $F$ determines the real complex structure
$J$ whose complexification acts by $i$ on $F$ and by $-i$ on
$\overline F$. Complex conjugation commutes with this
operator, so it restricts to the real space. The inverse map
takes the $+i$ eigenspace. Put $g=I+J^tJ$. Then $g$ is
positive definite and $J^tgJ=g$. With
\[
 Q_r=((1-r)I+rg)^{1/2},\qquad J_r=Q_rJQ_r^{-1},
                   \quad 0\le r\le1,
\]
all $\det Q_r$ are positive and
$J_1^tJ_1=g^{-1/2}J^tgJg^{-1/2}=I$.
An orthogonal $J$ has $g=2I$, so this deformation fixes it.
This is a strong deformation retraction onto the relevant
orthogonal component, represented by the conjugation orbit
$SO(6)/U(3)$.

The sphere fibrations $U(n-1)\to U(n)\to S^{2n-1}$ give
$\pi_1U(3)=\mathbb Z$. Starting with $SO(3)=\mathbb{RP}^3$,
the fibrations $SO(n-1)\to SO(n)\to S^{n-1}$ for
$n=4,5,6$ give $\pi_1SO(6)=\mathbb Z/2$.
The unitary loop rotating one complex coordinate maps to a
real two-plane rotation through $2\pi$. Its quaternion lift
in $SO(3)$ starts at $1$ and ends at $-1$, so it is the
nontrivial element; the later inclusions retain that class.
Thus $\pi_1U(3)\to\pi_1SO(6)$ is onto, and the homotopy
sequence gives $\pi_1(SO(6)/U(3))=0$. The fibre is also
connected. The map $P$ gives the same conclusions for $PD_3$.

Retain the local section on a slightly larger disc than $K$.
Take a locally finite triangulation of its exterior, relative
to the boundary. The period bundle is trivial over each cell.
Extend the section over vertices and edges using connectedness
of the fibre, and over two-cells using its zero fundamental
group. These extensions agree on the attaching boundaries;
local finiteness gives a continuous global section, still
holomorphic near $K$. The point $p_0$ remains inside the
retained local disc throughout this construction.

\item \textbf{The holomorphic section and actual period coordinates.}
Apply \ref{maps:src:F.3.10} to $(V_D,\mathscr U)$, with
$S=\{p_0\}$ and this Runge set $K$. The off-$S$ LSAP
bundle input is \ref{maps:src:F.3.8}; its section formulation
agrees with \cite[Definition~7.7 and Proposition~7.13(d)]{AlarconForstneric2025Sections}.
The lemma gives a holomorphic section over $B$ with the
prescribed finite jet at $p_0$. Its value there is in the
central graph chart, whose map $C\mapsto F_C$ is
biholomorphic. Its correction $C(t)$ is consequently
holomorphic at zero with exactly the requested matrix jet.
Undoing $T(s)$ gives \eqref{maps:eq:F.O6} and the original
$\rho$-equivariance on $B^*$.

For the local normalized period assertion, choose a local real
frame and let $\Lambda$ be its integer span in that chart; for an
integral representation retain the supplied lattice basis. Let
$\pi:\mathbb C^6\to E=\mathbb C^6/F$ be the actual quotient.
The images of the six real lattice basis vectors span $E$
over $\mathbb C$. Select three whose images are independent,
and let $P_\Lambda$ be the permutation matrix putting these
three first. In any local holomorphic frame of $E$, write
\[
 \pi P_\Lambda=[B,A],\qquad Z=B^{-1}A.
\]
The actual quotient coordinate change is left multiplication
by $B^{-1}$, and the normalized matrix is $[I_3,Z]$.
Its kernel, in this permuted lattice marking, is the graph
$\operatorname{im}\binom{-Z}{I_3}$. For real vectors $a,b$,
the real-coordinate map of the quotient is
\[
 \binom ab\longmapsto
 \begin{pmatrix}I_3&\operatorname{Re}Z\\0&\operatorname{Im}Z\end{pmatrix}
 \binom ab.
\]
It is invertible precisely when $\det\operatorname{Im}Z\ne0$;
its inverse sends $(r,s)$ to
$(r-\operatorname{Re}Z(\operatorname{Im}Z)^{-1}s,
 (\operatorname{Im}Z)^{-1}s)$.
The kernel graph satisfies
$\det(F,\overline F)=8i\det\operatorname{Im}Z$ in these
coordinates, consistent with its graph variable $-Z$.
The determinant factor of $P_\Lambda$ is retained by the
actual marking. Independence of the chosen three columns
persists locally, since $\det B$ remains nonzero.
This proves the asserted local period interpretation without
requiring one Grassmannian chart on the entire base.

\item \textbf{The specified survivor representations.}
For an F.2 survivor, \ref{maps:src:F.2.7} gives the genuine
two-point case, and \ref{maps:src:F.2.15} supplies a real
complex structure commuting with each finite matrix.
In a real basis induced by its complex basis, such an action
has matrix $R=\left(\begin{smallmatrix}A&-B\\B&A\end{smallmatrix}\right)$.
The exact complex matrix
\[
 C=\begin{pmatrix}I_3&I_3\\-iI_3&iI_3\end{pmatrix}
 \quad\hbox{satisfies}\quad
 C^{-1}RC=\operatorname{diag}(A+iB,A-iB).
\]
Hence its real determinant is $|\det(A+iB)|^2>0$.
Real conjugation back to the supplied lattice basis preserves
this determinant. The cusp has determinant one, so the
representation takes values in $GL^+(6,\mathbb R)$.
The integral cusp marking is supplied by the first item.
All the preceding constructions therefore apply to each
specified survivor and give the literal period family on
$B^*$. The other points remain punctured throughout.
\end{enumerate}
\end{proof}
